% !TeX root = ../main.tex

\begin{translation}
\label{cha:translation}

\title{书面翻译题目}
\maketitle

\tableofcontents


\section{单目标规划}
北冥有鱼，其名为鲲。鲲之大，不知其几千里也。化而为鸟，其名为鹏。鹏之背，不知其几
千里也。怒而飞，其翼若垂天之云。是鸟也，海运则将徙于南冥。南冥者，天池也。\cite{abrahams99tex}
\begin{equation}\tag*{(123)}
 p(y|\mathbf{x}) = \frac{p(\mathbf{x},y)}{p(\mathbf{x})}=
\frac{p(\mathbf{x}|y)p(y)}{p(\mathbf{x})}
\end{equation}

吾生也有涯，而知也无涯。以有涯随无涯，殆已！已而为知者，殆而已矣！为善无近名，为
恶无近刑，缘督以为经，可以保身，可以全生，可以养亲，可以尽年。

\subsection{线性规划}
庖丁为文惠君解牛，手之所触，肩之所倚，足之所履，膝之所倚，砉然响然，奏刀騞然，莫
不中音，合于桑林之舞，乃中经首之会。
\begin{table}[ht]
\centering
  \centering
  \caption{这是一个表格例子}
  \label{tab:badtabular3}
  \begin{tabular}[c]{|m{1.5cm}|c|c|c|c|c|c|}\hline
    \multicolumn{2}{|c|}{Network Topology} & \# of nodes &
    \multicolumn{3}{c|}{\# of clients} & Server \\\hline
    GT-ITM & Waxman Transit-Stub & 600 &
    \multirow{2}{2em}{2\%}&
    \multirow{2}{2em}{10\%}&
    \multirow{2}{2em}{50\%}&
    \multirow{2}{1.2in}{Max. Connectivity}\\\cline{1-3}
    \multicolumn{2}{|c|}{Inet-2.1} & 6000 & & & &\\\hline
    \multirow{2}{1.5cm}{Xue} & Rui  & Ni &\multicolumn{4}{c|}{\multirow{2}*{\thuthesis}}\\\cline{2-3}
    & \multicolumn{2}{c|}{ABCDEF} &\multicolumn{4}{c|}{} \\\hline
\end{tabular}
\end{table}

文惠君曰：“嘻，善哉！技盖至此乎？”庖丁释刀对曰：“臣之所好者道也，进乎技矣。始臣之
解牛之时，所见无非全牛者；三年之后，未尝见全牛也；方今之时，臣以神遇而不以目视，
官知止而神欲行。依乎天理，批大郤，导大窾，因其固然。技经肯綮之未尝，而况大坬乎！
良庖岁更刀，割也；族庖月更刀，折也；今臣之刀十九年矣，所解数千牛矣，而刀刃若新发
于硎。彼节者有间而刀刃者无厚，以无厚入有间，恢恢乎其于游刃必有余地矣。是以十九年
而刀刃若新发于硎。虽然，每至于族，吾见其难为，怵然为戒，视为止，行为迟，动刀甚微，
謋然已解，如土委地。提刀而立，为之而四顾，为之踌躇满志，善刀而藏之。”

文惠君曰：“善哉！吾闻庖丁之言，得养生焉。”


\subsection{非线性规划}
孔子与柳下季为友，柳下季之弟名曰盗跖。盗跖从卒九千人，横行天下，侵暴诸侯。穴室枢
户，驱人牛马，取人妇女。贪得忘亲，不顾父母兄弟，不祭先祖。所过之邑，大国守城，小
国入保，万民苦之。孔子谓柳下季曰：“夫为人父者，必能诏其子；为人兄者，必能教其弟。
若父不能诏其子，兄不能教其弟，则无贵父子兄弟之亲矣。今先生，世之才士也，弟为盗
跖，为天下害，而弗能教也，丘窃为先生羞之。丘请为先生往说之。”
\begin{figure}[h]
  \centering
  \includegraphics{thu-whole-logo.pdf}
  \caption{这是图片的例子}
  \label{tab:badfigure3}
\end{figure}

柳下季曰：“先生言为人父者必能诏其子，为人兄者必能教其弟，若子不听父之诏，弟不受
兄之教，虽今先生之辩，将奈之何哉？且跖之为人也，心如涌泉，意如飘风，强足以距敌，
辩足以饰非。顺其心则喜，逆其心则怒，易辱人以言。先生必无往。”

孔子不听，颜回为驭，子贡为右，往见盗跖。

\subsection{整数规划}
盗跖乃方休卒徒大山之阳，脍人肝而餔之。孔子下车而前，见谒者曰：“鲁人孔丘，闻将军
高义，敬再拜谒者。”谒者入通。盗跖闻之大怒，目如明星，发上指冠，曰：“此夫鲁国之
巧伪人孔丘非邪？为我告之：尔作言造语，妄称文、武，冠枝木之冠，带死牛之胁，多辞缪
说，不耕而食，不织而衣，摇唇鼓舌，擅生是非，以迷天下之主，使天下学士不反其本，妄
作孝弟，而侥幸于封侯富贵者也。子之罪大极重，疾走归！不然，我将以子肝益昼餔之膳。”


% 书面翻译的参考文献
\bibliographystyle{plainnat}
\bibliography{ref/appendix}

% 也可以使用 thebiliography 环境手写
% \begin{thebibliography}{2}
%   \bibitem{abrahams99tex}
%   P.~W. Abrahams, K.~Berry, and K.~A. Hargreaves, \emph{{\TeX} for the
%     Impatient}.\hskip 1em plus 0.5em minus 0.4em\relax Addison-Wesley, 1990.

%   \bibitem{salomon1995advanced}
%   D.~Salomon, ``The advanced {\TeX}book.''\hskip 1em plus 0.5em minus 0.4em\relax
%     New York: Springer, 1995.
% \end{thebibliography}

% 书面翻译对应的原文索引
\begin{translation-index}
  \nocite{salomon1995advanced}
  \bibliographystyle{plainnat}
  \bibliography{ref/appendix}
\end{translation-index}

\end{translation}
